% ID: SEL-D0005-OA-E04
% ESTADO: cotejada. Segunda pasada visual completa y borrador compilado revisados.
% Fuente: Selectividad/sel_2019_matematicas_aplicadas/Reserva A/Exámenes Mat Aplicadas CCSS - A y B.pdf
% SHA-256: d28871f7f203cf1be9e479e9d642cdb6b1623704a8b76a2b57cfa2d12950f293
% Página: 1; opción A; ejercicio 4.

Para estimar la proporción de empleados de una empresa que usan lentillas,
se toma una muestra al azar de 60 empleados de la misma y se observa que
16 usan lentillas.
\begin{enumerate}[label=\alph*)]
  \item \textbf{(1.5 puntos)} Halle, con un nivel de confianza del 90\,\%,
  un intervalo para estimar la proporción.
  \item \textbf{(1 punto)} Con el mismo nivel de confianza del apartado
  anterior y manteniendo la misma proporción muestral, determine el tamaño
  mínimo de la muestra para que el error cometido en la estimación de la
  proporción sea inferior a $0.1$.
\end{enumerate}
